\exercisesheader{}

% 37

\eoce{\qt{Berat kucing\label{cat_weights}} Histogram di bawah ini merepresentasikan
berat (dalam kg) 47 kucing betina dan 97 kucing jantan. \footfullcite{cats} \\
\begin{minipage}[c]{0.47\textwidth}
\begin{parts}
\item Berapa proporsi kucing-kucing ini yang beratnya kurang dari 2.5 kg?
\item Berapa proporsi kucing-kucing ini yang beratnya antara 2.5 dan 2.75 kg?
\item Berapa proporsi kucing-kucing ini yang beratnya antara 2.75 dan 3.5 kg?
\end{parts} \vspace{27mm}
\end{minipage}
\begin{minipage}[c]{0.05\textwidth}
$\:$
\end{minipage}
\begin{minipage}[c]{0.48\textwidth}
\begin{center}
\Figures[Ditampilkan histogram berat badan kucing dalam kilogram. Rentang beratnya dari 2.0 sampai 4.0 dan setiap bin histogram memiliki lebar 0.25. Tinggi kedelapan bin, dari kiri ke kanan, adalah 29, 32, 21, 26, 12, 15, 5, dan 4.]{}{eoce/cat_weights}{cat_weights}
\end{center}
\end{minipage}
}{}

% 38

\eoce{\qt{Pendapatan dan gender\label{income_gender}} Tabel frekuensi relatif
di bawah ini menampilkan distribusi pendapatan pribadi total tahunan (dalam dolar
yang disesuaikan dengan inflasi tahun 2009) untuk sampel representatif yang terdiri atas
96,420,486 penduduk Amerika Serikat. Data ini berasal dari American Community Survey
tahun 2005--2009. Sampel ini terdiri atas 59\% laki-laki dan 41\% perempuan. \footfullcite{acsIncome2005-2009} \\

\noindent\begin{minipage}[c]{0.60\textwidth}
\begin{parts}
\item Deskripsikan distribusi pendapatan pribadi total.
\item Berapa peluang bahwa penduduk AS yang dipilih secara acak berpenghasilan kurang dari
\$50,000 per tahun?
\item Berapa peluang bahwa penduduk AS yang dipilih secara acak berpenghasilan kurang dari
\$50,000 per tahun dan berjenis kelamin perempuan? Nyatakan setiap asumsi yang Anda buat.
\item Sumber data yang sama menunjukkan bahwa 71.8\% perempuan berpenghasilan kurang dari
\$50,000 per tahun. Gunakan nilai ini untuk menentukan apakah asumsi yang Anda buat pada bagian (c) valid.
\end{parts}
\end{minipage}
\begin{minipage}[c]{0.4\textwidth}
{\small
\begin{center}
\begin{tabular}{lr}
  \hline
\textit{Pendapatan}         & \textit{Total} \\
  \hline
\$1 sampai \$9,999 atau rugi & 2.2\% \\
\$10,000 sampai \$14,999    & 4.7\% \\
\$15,000 sampai \$24,999    & 15.8\% \\
\$25,000 sampai \$34,999    & 18.3\% \\
\$35,000 sampai \$49,999    & 21.2\% \\
\$50,000 sampai \$64,999    & 13.9\% \\
\$65,000 sampai \$74,999    & 5.8\% \\
\$75,000 sampai \$99,999    & 8.4\% \\
\$100,000 atau lebih        & 9.7\% \\
   \hline
\end{tabular}
\end{center}
}
\end{minipage}
}{}
